\section{What the mapper asked for in return}
\label{sec:ch04-what-the-mapper-asked-for}
\index{Entity Framework Core!model binding|(}

Mosaic's domain types are records with positional constructors, and they were
written that way in chapter~\ref{ch:hotchocolate-quickly} because immutability is
the right default for something a resolver hands out. Entity Framework Core
disagreed with three of them. The disagreements are small, the reason behind all
three is the same, and the reason is not quite what the error message says.

\index{Entity Framework Core!collection navigations}%
The first complaint was easy and I fixed it without thinking. A collection
navigation has to be something the mapper can add to while it materialises a
row, and the compiler-generated backing field of a record property has the
declared type, so an \code{IReadOnlyList} leaves it nothing to fill.
\code{Order.Lines} became a \code{List<OrderLine>}. That is one interface on one
property, and it does not change the GraphQL field at all.

Then, with the collection type already fixed, the model still would not build:

\begin{minted}[fontsize=\footnotesize]{text}
System.InvalidOperationException: No suitable constructor was found for the
type 'Order'. The following constructors had parameters that could not be
bound to properties of the type:
    Cannot bind 'Lines' in 'Order(Guid Id, Guid CustomerId,
        DateTimeOffset PlacedAt, List<OrderLine> Lines)'
Note that only mapped properties can be bound to constructor parameters.
Navigations to related entities, including references to owned types, cannot
be bound.
\end{minted}

\index{Entity Framework Core!constructor binding}%
Clear, specific, and it arrives at startup rather than on the first query, which
is chapter~\ref{ch:hotchocolate-quickly}'s eager initialisation paying out
again. The documented rule is that a constructor parameter binds to a mapped
property and a navigation is not one~\autocite{microsoft2026efconstructors}. An
owned type is an entity without a key of its own, so it arrives as a navigation,
so it cannot be a constructor parameter. Fine.

\subsection*{The half of the rule that is not written down}

\index{Entity Framework Core!complex types}%
\code{Money} is a value. It has an amount and a currency, no identity, and no
reason to exist apart from the thing it prices. Entity Framework Core has a
feature for exactly that, distinct from owned types: a complex property, which
maps to columns on the same row and is treated as a value rather than as a
tracked entity. Since a complex property is a property and not a navigation, the
error message above appeared to be telling me how to fix it.

It was not. Moving \code{ProductPrice.Amount} from \code{OwnsOne} to
\code{ComplexProperty} produced the identical exception, with \code{Amount}
named where \code{Lines} had been. Against Entity Framework Core 10.0.10, only
scalars bind to constructor parameters. Neither an owned reference nor a complex
property qualifies, and the message names only the first of those.

I am stating that as a measurement rather than as a rule, because it is one. I
compiled it, ran it, and read the exception; I did not find it on a page. If a
later release relaxes it, this paragraph is the thing that dates, and a reader
on a newer version should try the positional constructor before believing me.

\index{Entity Framework Core!init properties}%
The fix is unexciting and identical in all three places. The value leaves the
primary constructor and becomes an \code{init} property, which Entity Framework
Core populates through the backing field:

\begin{minted}{csharp}
public sealed record ProductPrice(Guid ProductId)
{
    public required Money Amount { get; init; }
}
\end{minted}

Three properties moved out of three constructors. Together with the collection
type above, that is every concession Mosaic made to the mapper, and none of them
reaches the graph.

\subsection*{A key nobody asked for}

\index{Entity Framework Core!shadow properties}%
With \code{Money} settled as a complex property, order lines could not stay
owned, because an owned collection is a navigation and \code{OrderLine} is a
record whose price wants to be a complex property of its own. So the line became
a related entity. It has no identifier in the domain model and none on the
graph, and a relational row needs one, so both the key and the foreign key back
to the order live as shadow properties:

\begin{minted}{csharp}
builder.ToTable("order_lines");

builder.Property<int>("Id").ValueGeneratedOnAdd();
builder.HasKey("Id");

builder.ComplexProperty(l => l.UnitPrice, money =>
{
    money.Property(m => m.Amount).HasPrecision(12, 2).IsRequired();
    money.Property(m => m.Currency).HasMaxLength(3).IsRequired();
});
\end{minted}

A shadow property exists in the database and in the mapper's model, and nowhere
in the C\#. No code can read it or set it. When the thing you are modelling
genuinely has no identity worth exposing, this is better than inventing an
\code{Id} property and then remembering to hide it from everything.

\subsection*{Making the SQL readable}

\index{PostgreSQL!identifier folding}%
One more piece of mapping earns its place, because the rest of this chapter is
spent reading generated statements. PostgreSQL folds unquoted identifiers to
lower case, so a table Entity Framework Core calls \code{StockLevels} has to be
quoted every time anybody writes it by hand. Renaming everything once at model
build removes that:

\begin{minted}[fontsize=\footnotesize]{csharp}
foreach (var entity in modelBuilder.Model.GetEntityTypes())
{
    if (entity.GetTableName() is { } table)
    {
        entity.SetTableName(ToSnakeCase(table));
    }

    foreach (var property in entity.GetProperties())
    {
        property.SetColumnName(ToSnakeCase(property.GetColumnName()));
    }
    ...
}
\end{minted}

\code{GetProperties} does not return complex ones, and I found that out the way
you would expect. The price columns turned up quoted and in mixed case, in the
middle of an otherwise lower-case statement I was about to put in this chapter:

\begin{minted}[fontsize=\footnotesize]{sql}
SELECT p.product_id, p."Amount_Amount", p."Amount_Currency"
FROM prices AS p
WHERE p.product_id = ANY (@productIds)
\end{minted}

Complex properties need a second pass, over \code{GetComplexProperties}.

\subsection*{And the schema did not move}

\index{schema!stability}%
Here is the part worth stopping on. Everything above changed: the storage, the
mapping, three type definitions, the service constructors, six domain services
rewritten against \code{IQueryable}. The exported SDL at tag \code{ch04-ef} is
byte-identical to the one at tag \code{ch03}.

That is a consequence of a decision made two chapters ago, and it is the one
that pays for itself first. Chapter~\ref{ch:hotchocolate-quickly}
put the \code{price} field in Pricing's folder rather than on Catalog's
\code{Product} record, and the reason it did is the same reason the whole data
layer could be replaced without a client noticing. A GraphQL
schema is a contract about fields and types, and where a field's value comes
from is not part of it. A resolver is where a change of mind can go without a
client noticing.

Hold on to that, because the next section changes every one of those resolvers
and the schema does not move then either.
\index{Entity Framework Core!model binding|)}
